
\chapter{Processi stocastici a tempo discreto}

Introduciamo il concetto di processo stocastico a tempo discreto. I processi stocastici sono oggetti matematici il cui scopo è descrivere l'evoluzione temporale di fenomeni aleatori. 
Senza perdita di generalità, l'indice temporale sarà identificato con un sottoinsieme dei numeri naturali, solitamente a partire da $n=0$. Indicheremo l'insieme dei tempi con $\N$ (orizzonte finito) oppure $\N_0$ (orizzonte infinito). Nelle formule useremo $\N_0 = \N \cup \{0\}$.

\section{Nozione di Processo Stocastico}

Assumiamo fissato uno spazio di probabilità $(\Omega, \F, \PP)$ e uno spazio di arrivo $E$ (al più numerabile), munito della $\sigma$-algebra discreta (l'insieme delle parti).

\begin{definition}[Processo Stocastico]
Un processo stocastico a valori in $E$ è una famiglia $X = (X_n)_{n \in \N_0}$ di variabili aleatorie $X_n : \Omega \to E$.
\end{definition}

Possiamo pensare a $X$ come alla realizzazione in sequenza di variabili aleatorie nel tempo. Se fissiamo un esito $\omega \in \Omega$, la sequenza numerica (o a valori in $E$):
\[ (X_0(\omega), X_1(\omega), X_2(\omega), \dots) \]
prende il nome di \emph{traiettoria} del processo stocastico $X$.

\section{Filtrazioni e Struttura dell'Informazione}

Vogliamo dotare il modello di una struttura di informazione che rifletta l'apprendimento progressivo nel tempo: man mano che il tempo scorre, l'individuo "legge" cosa succede.
Supponiamo che l'individuo osservi i valori di un processo $\xi = (\xi_n)$. Al tempo $n$, egli conosce i valori $\xi_0(\omega), \dots, \xi_n(\omega)$. Questa conoscenza pregressa può modificare le aspettative sul futuro.
Formalizziamo questa idea con le \emph{filtrazioni}.

\begin{definition}[Filtrazione Naturale]
Sia $\xi$ un processo stocastico. Definiamo:
\[ \F_n^\xi = \sigma(\xi_0, \xi_1, \dots, \xi_n) \]
cioè la più piccola $\sigma$-algebra rispetto a cui tutte le variabili $\xi_0, \dots, \xi_n$ sono misurabili.
Chiaramente $\F_n^\xi \subseteq \F_{n+1}^\xi$. La famiglia crescente $(\F_n^\xi)_{n \in \N_0}$ si chiama \emph{filtrazione generata} da $\xi$ o \emph{filtrazione naturale}.
\end{definition}

\begin{definition}[Processo Adattato]
Un processo stocastico $X = (X_n)$ si dice \emph{adattato} alla filtrazione $(\F_n)$ se, per ogni tempo $n$, la variabile aleatoria $X_n$ è misurabile rispetto a $\F_n$.
\end{definition}

In parole povere, $X$ è adattato se il suo valore al tempo $n$ dipende unicamente dall'informazione disponibile fino al tempo $n$, senza sbirciare nel futuro. 
Spesso si usa la filtrazione naturale generata da $X$ stesso: in tal caso, il processo è banalmente adattato.

\section{Tempi di Arresto (Stopping Times)}

Una nozione cruciale per lo studio dei processi stocastici è quella di \emph{tempo di arresto} (stopping time). Formalizza l'idea di una regola d'arresto che si basa unicamente sull'informazione presente e passata, mai su quella futura.

\begin{definition}[Tempo di arresto]
Sia $(\F_n)_{n \in \N_0}$ una filtrazione. Una variabile aleatoria $\tau : \Omega \to \N_0 \cup \{\infty\}$ è detta \emph{tempo di arresto} se, per ogni $n \in \N_0$:
\[ \{\tau \le n\} \in \F_n \]
\end{definition}

\begin{hint}[Condizione equivalente]
Dire che $\{\tau \le n\} \in \F_n$ equivale a dire $\{\tau = n\} \in \F_n$.
A parole: l'informazione disponibile al tempo $n$ è sufficiente per stabilire inequivocabilmente se l'evento in questione "è già accaduto" o no.
\end{hint}

Un classico esempio di tempo di arresto è il \emph{tempo di primo passaggio} (hitting time).

\begin{definition}[Tempo di primo passaggio]
Sia $X$ un processo adattato e sia $A \subseteq E$. Il tempo di primo ingresso (o hitting time) in $A$ è definito come:
\[ \tau_A = \inf\{n \in \N_0 : X_n \in A\} \]
(con la convenzione $\inf \emptyset = \infty$).
\end{definition}

\begin{esempio}[Cosa \textbf{non} è un tempo di arresto]
Il tempo in cui un processo su orizzonte limitato $N$ raggiunge il suo \emph{massimo globale}, cioè $\tau = \operatorname{argmax}_{k \le N} X_k$, \textbf{non} è in generale un tempo di arresto. Per stabilire se $\tau = n$, bisognerebbe sapere se nei tempi futuri $n+1, \dots, N$ ci sarà un valore ancora più alto, ma questo richiede di "sbirciare nel futuro" (non è in $\F_n$).
\end{esempio}

\begin{definition}[Processo arrestato]
Dato un processo adattato $X_n$ e un tempo di arresto $\tau$, il \emph{processo arrestato} (stopped process) a $\tau$ è definito da:
\[ X_n^\tau = X_{\min(n, \tau)} = \begin{cases}
X_n & \text{se } n < \tau \\
X_\tau & \text{se } n \ge \tau
\end{cases} \]
\end{definition}

\begin{theorem}[Il processo arrestato è adattato]
Se $X$ è adattato alla filtrazione $(\F_n)$ e $\tau$ è un tempo di arresto, allora anche il processo arrestato $X^\tau$ è adattato alla medesima filtrazione.
\end{theorem}

\section{Esercizi Selezionati (Capitolo 3)}

\subsection*{Esercizio 1: Minimo di tempi di arresto}
Siano $\tau_1$ e $\tau_2$ due tempi di arresto. Dimostrare che $\tau = \min(\tau_1, \tau_2)$ è anch'esso un tempo di arresto.
\begin{spiegazione}
Per definizione, vogliamo verificare se $\{\tau \le n\} \in \F_n$. Osserviamo che il minimo tra due numeri è $\le n$ se e solo se almeno uno dei due è $\le n$:
\[ \{\min(\tau_1, \tau_2) \le n\} = \{\tau_1 \le n\} \cup \{\tau_2 \le n\} \]
Poiché entrambi gli eventi appartengono a $\F_n$ e $\F_n$ è una $\sigma$-algebra (chiusa rispetto alle unioni), l'unione appartiene a $\F_n$. Quindi $\tau$ è un tempo di arresto.
\end{spiegazione}

\subsection*{Esercizio 2: Tempo di ultima uscita (Last exit time)}
Sia $L_A = \sup\{n \in \N_0 : X_n \in A\}$ il tempo dell'ultima visita di un processo all'insieme $A$. Perché in generale \emph{non} è un tempo di arresto?
\begin{spiegazione}
Per determinare se $\{L_A = n\}$ si è verificato, non basta sapere che $X_n \in A$, ma occorre anche avere la certezza che per ogni $k > n$ il processo non tornerà in $A$ ($X_k \notin A$). Ma questa è un'informazione sul futuro, non contenuta in $\F_n$. Dunque non è una regola d'arresto (non puoi "fermarti" a $n$ sapendo che è l'ultima volta che tocchi $A$, perché non lo saprai finché non finirà il tempo!).
\end{spiegazione}

\subsection*{Esercizio 3: Atomi di una filtrazione}
Se $\xi_n \sim B(p)$ (lancio di moneta i.i.d.), descrivere gli atomi della $\sigma$-algebra $\F_1 = \sigma(\xi_0, \xi_1)$.
\begin{spiegazione}
Gli atomi (eventi elementari minimali) sono tutte le possibili realizzazioni congiunte delle prime due variabili, che descrivono completamente lo stato di conoscenza a tempo $n=1$. Poiché ciascuna può assumere 2 valori, vi sono $2^2 = 4$ atomi:
$A_1 = \{\xi_0=0, \xi_1=0\}$, $A_2 = \{\xi_0=0, \xi_1=1\}$, $A_3 = \{\xi_0=1, \xi_1=0\}$, $A_4 = \{\xi_0=1, \xi_1=1\}$.
Qualsiasi evento misurabile $\F_1$ è unione di questi atomi.
\end{spiegazione}

\subsection*{Esercizio 4: Somma di tempi di arresto}
Siano $\tau_1, \tau_2$ tempi di arresto. Mostrare che $\tau_1 + \tau_2$ è un tempo di arresto.
\begin{spiegazione}
L'evento $\{\tau_1 + \tau_2 = n\}$ si può scomporre come unione di tutti i modi in cui i due tempi possono sommare a $n$:
\[ \{\tau_1 + \tau_2 = n\} = \bigcup_{k=0}^n \big(\{\tau_1 = k\} \cap \{\tau_2 = n-k\}\big) \]
Per $k \le n$, l'evento $\{\tau_1 = k\} \in \F_k \subseteq \F_n$. Anche l'evento per $\tau_2$ appartiene a $\F_{n-k} \subseteq \F_n$. L'intersezione e la loro unione finita sono in $\F_n$. Dunque la somma è un tempo di arresto.
\end{spiegazione}
